\begin{proposition}[the mapping-torus comparison]\label{push:src:C.1.1}
For an orientation-preserving smooth diffeomorphism
$\phi:S^6\to S^6$, put
\[
 M_\phi=(S^6\times[0,1])/((x,1)\sim(\phi(x),0)),\qquad
 \Sigma_\phi=D^7_+\cup_\phi(-D^7_-).
\]
The boundary identification uses the usual oriented parametrizations
of the two underlying copies of $S^6$; the minus sign reverses the
orientation of the second disc. Then
\[
 M_\phi\cong (S^1\times S^6)\mathbin\#\Sigma_\phi
\]
as oriented smooth manifolds. With the corresponding sphere-gluing
identification, $[\phi]\mapsto[\Sigma_\phi]$ identifies
$\pi_0\Diff^+(S^6)$ with $\Theta_7$. For either spin structure on
$M_\phi$,
\[
 \mu(M_\phi)=\mu(\Sigma_\phi).
\]
For the family in the prompt the gluing map is
$\Phi=F_{2,c_0-2}\circ H$, with $c(r)=c_0-2r$.
\end{proposition}

\begin{proof}
An ambient isotopy first arranges that $\phi$ is the identity on a
small closed disc: move the image of its centre back, correct the
oriented derivative through $\GL^+(6,\R)$, and use isotopy extension
on a smaller coordinate disc. Consequently its support lies in the
complementary disc $D^6$. Cut the mapping torus at its seam and take
the block $D^6\times[-\epsilon,\epsilon]$. The lower half of this
block has the old coordinates; on the upper half replace the
coordinates by $\phi^{-1}$. This identifies the reglued block with a
standard seven-disc. Its attaching map is $\phi$ on its upper
hemisphere and the identity on the lower hemisphere and side collar.
The upper hemisphere has boundary orientation $o_{D^6}$ for the
orientation $dt\wedge o_{D^6}$ on the block. Thus the replacement
has exactly the sign in $D^7_+\cup_\phi(-D^7_-)$.
Removing a small ball from the second disc in that sphere leaves the
same attaching construction; this proves the connected-sum formula.

Sphere gluing identifies pseudo-isotopy classes with homotopy spheres;
Cerf's pseudo-isotopy theorem changes pseudo-isotopy to isotopy for
the simply connected six-sphere. The exact primary statements are
Theorem 0 and Corollaries 1--2 of \citet[p.~8]{Cerf1970};
the sphere case is also explained in Theorem 1 and Corollary 1 of
\citet[p.~156]{Poenaru1971}. These statements compare the
sphere gluing and extension-over-a-disc constructions. This is the
high-dimensional theorem, not the separate statement $\Gamma_4=0$.
The oriented connected-sum group is Theorem 1.1 and Lemma 2.1 of
\citet[pp.~504--505]{KervaireMilnor1963}; their dimension table on
p.~504 gives order 28 in dimension seven.

The sphere summand is spin. The product $S^1\times S^6$ is spin and
has two spin structures; both bound $S^1\times D^7$ with the
corresponding choice on $S^1$. That coboundary has zero fourth
cohomology and zero signature, hence zero Eells--Kuiper invariant.
Boundary connected sum of this coboundary with a spin coboundary of
$\Sigma_\phi$ proves additivity directly: the fourth intersection
forms are an orthogonal sum and the Pontryagin evaluations add.
Thus the two boundary spin structures give the same value. This is
also the connected-sum property in \citet[Section 3]{EellsKuiper1962}.

Finally $e^{-\pi i(c_0-2r)}=e^{-\pi ic_0}e^{2\pi ir}$.
Parallel transport therefore reaches $X_{c_0-2}$ along the positive
$w$-circle; its endpoint is identified by $F_{2,c_0-2}$.
This confirms the sign and endpoint in \cref{fam:src:E.2.2}. Two smooth connections
can be joined by a smooth path of horizontal distributions, producing
an isotopy of the resulting monodromies. Holomorphic automorphisms in
the connected group $\C^*$ also change the result only by isotopy.
\end{proof}

